\documentclass[11pt]{article}
\usepackage[margin=1in]{geometry}
\usepackage{amsmath,amssymb,bm}
\usepackage{graphicx}
\usepackage{booktabs}
\usepackage{xcolor}
\usepackage[hidelinks]{hyperref}
\usepackage{listings}
\usepackage{enumitem}

\lstset{basicstyle=\ttfamily\small,breaklines=true,frame=single,
  backgroundcolor=\color{black!4},columns=fullflexible,showstringspaces=false}
\newcommand{\xv}{\bm{x}}
\newcommand{\uv}{\bm{u}}
\newcommand{\xiv}{\bm{\xi}}
\newcommand{\Ften}{\bm{F}}
\newcommand{\bten}{\bm{b}}
\newcommand{\grad}{\nabla}
\newcommand{\dd}[2]{\frac{\partial #1}{\partial #2}}

\title{\textbf{pyRMT --- Technical Manual}\\[2mm]
\large A staggered-grid Reference Map Technique for 2D incompressible
fluid--structure interaction}
\author{Saman Seifi \\ \small Method, discretization, and implementation reference}
\date{\today}

\begin{document}
\maketitle
\begin{abstract}
\noindent
pyRMT is a research code for two-dimensional, incompressible fluid--structure
interaction (FSI) in which a hyperelastic (or viscoelastic) solid and a Newtonian
fluid share a single velocity field, pressure, and fixed Eulerian grid. The solid
is tracked by the \emph{Reference Map Technique} (RMT): an Eulerian field
$\xiv(\xv,t)$ that records, for each spatial point currently occupied by solid,
its original (material) position. This manual documents the governing equations,
the constitutive models (neo-Hookean and finite-strain viscoelasticity via the
log-conformation method), the spatial and temporal discretization on a
Marker-and-Cell (MAC) staggered grid with an \emph{exact} discrete projection, the
interface, surface-tension and contact treatments, and the verification and
validation evidence. A companion \emph{User Manual} covers installation, the public
API, and how to run the benchmarks.
\end{abstract}

\tableofcontents

\section{Scope and two solver families}
The package contains two solver families that share the same RMT idea but differ in
their grid arrangement:
\begin{itemize}[nosep]
  \item \textbf{Collocated} (\texttt{pyRMT/functions.py}): velocity and pressure at
  cell centres, Rhie--Chow interpolation to suppress checkerboarding; direct DCT/FFT
  Poisson solves and an AMG fallback for variable density. This is the original
  solver (after Jain, Kamrin \& Mani, 2019).
  \item \textbf{Staggered / MAC} (\texttt{pyRMT/mac.py}): pressure at cell centres,
  normal velocities on cell faces. The discrete divergence and gradient are exact
  transposes, so the projection is divergence-free to machine precision and no
  Rhie--Chow correction is needed. This is the recommended path for surface tension
  and contact, and is the basis for the viscoelastic extension.
\end{itemize}
Unless stated otherwise, this manual documents the \textbf{staggered/MAC} solver.

\section{Governing equations}
The one-fluid incompressible momentum and continuity equations on the whole domain
$\Omega$ are
\begin{align}
\rho\!\left(\dd{\uv}{t}+\uv\!\cdot\!\grad\uv\right)
&= -\grad p + \grad\!\cdot\!\bm{\sigma}_{\!d} + \bm{f}, \\
\grad\!\cdot\!\uv &= 0,
\end{align}
where $p$ is the pressure (the incompressibility Lagrange multiplier),
$\bm{\sigma}_{\!d}$ is the deviatoric stress, and $\bm{f}$ collects body and
interfacial forces. The deviatoric stress is \emph{blended} between fluid and solid
through a smoothed indicator $H$ (Section~\ref{sec:interface}):
\begin{equation}
\bm{\sigma}_{\!d} = H\,\bm{\sigma}_{\!f} + (1-H)\,\bm{\sigma}_{\!s},
\qquad
\bm{\sigma}_{\!f} = 2\mu_f\,\bm{D}(\uv),
\label{eq:blend}
\end{equation}
with $\bm{D}=\tfrac12(\grad\uv+\grad\uv^{\!\top})$. Here $H=1$ in fluid and $H=0$
in solid. A single density $\rho$ is used by default (variable density
$\rho_s\neq\rho_f$ is available only in the collocated AMG path).

\section{The Reference Map Technique}
The reference map $\xiv(\xv,t)$ satisfies pure advection in the solid,
\begin{equation}
\dd{\xiv}{t}+\uv\!\cdot\!\grad\xiv = \bm{0},
\qquad \xiv(\xv,0)=\xv ,
\label{eq:xi}
\end{equation}
so $\xiv$ maps a current point back to its material point. The deformation gradient
and left Cauchy--Green (Finger) tensor follow from $\xiv$:
\begin{equation}
\Ften = (\grad\xiv)^{-1},\qquad \bten = \Ften\Ften^{\!\top},\qquad J=\det\Ften .
\end{equation}
The solid region is recovered from a reference (material) level set $\phi_0$ via the
compatibility relation $\phi(\xv,t)=\phi_0(\xiv(\xv,t))$; the zero contour is the
fluid--solid interface. Outside the solid the reference map is undefined, so after
each advection step it is \emph{extrapolated} a few cells into the surrounding band
(Section~\ref{sec:interface}).

\section{Constitutive models}
\subsection{Incompressible neo-Hookean solid}
The default solid stress is the neo-Hookean Cauchy stress in the convention of
Rycroft \emph{et al.} (2018),
\begin{equation}
\bm{\sigma}_{\!s} = \mu_s\,\bten ,
\label{eq:neo}
\end{equation}
with shear modulus $\mu_s$. The hydrostatic part is carried by the pressure (the
material is incompressible, $\kappa=0$ by default), so only the deviatoric part of
\eqref{eq:neo} is dynamically relevant.
\paragraph{Convention note.} Some references (e.g.\ Jain \emph{et al.}, 2019) write
$\bm\sigma_s = 2\mu_s\,\mathrm{dev}(\bten)$. The factor difference is purely a
definition of $\mu_s$ (the strain energy is written with $\tfrac12$ vs.\ $1$); it is
\emph{not} a discrepancy. An optional volume-corrected (isochoric) form
$\bm\sigma_s=\mu_s J^{-2}(\bten-\tfrac12\mathrm{tr}(\bten)\bm I)+\kappa(J-1)\bm I$ is
implemented in \texttt{solid\_cauchy\_stress} but is neutral on convergence and is
not the default.

\subsection{Finite-strain viscoelasticity (log-conformation)}
\label{sec:ve}
For viscoelastic solids the stress is carried by an \emph{elastic} Finger tensor
$\bten_e$ evolved by the upper-convected Maxwell (UCM) model with relaxation time
$\tau$,
\begin{equation}
\dd{\bten_e}{t}+(\uv\!\cdot\!\grad)\bten_e - \bm{L}\bten_e - \bten_e\bm{L}^{\!\top}
= -\frac{1}{\tau}\bigl(\bten_e-\bm I\bigr),
\qquad \bm L=\grad\uv ,
\label{eq:ucm}
\end{equation}
and the stress
\begin{equation}
\bm{\sigma}_{\!s} = G\,\bten_e \;+\; G_\infty\,\mathrm{dev}(\Ften\Ften^{\!\top}).
\label{eq:vestress}
\end{equation}
The optional second term is a non-relaxing equilibrium (Zener / standard-linear-solid)
branch, so that $G_\infty>0$ keeps the body a \emph{solid} (returns to shape) while
$G_\infty=0$ gives a Maxwell material (flows under sustained load). The limit
$\tau\!\to\!\infty$ recovers the neo-Hookean solid \eqref{eq:neo}.

\paragraph{Log-conformation.} Advecting $\bten_e$ directly fails on a grid: a
diffusive scheme artificially relaxes it (kills the elastic strain), while a
non-dissipative scheme loses positive-definiteness and blows up. Following Fattal \&
Kupferman (2004) we evolve $\bm\psi=\log\bten_e$, so $\bten_e=\exp(\bm\psi)$ is
symmetric positive-definite \emph{by construction} for any $\bm\psi$. In the
eigenframe of $\bm\psi$ the homogeneous evolution is
\begin{equation}
\dot{\bm\psi} = (\bm\Omega\bm\psi-\bm\psi\bm\Omega) + 2\bm B
 + \tfrac{1}{\tau}\bigl(e^{-\bm\psi}-\bm I\bigr),
\end{equation}
with the extensional ($\bm B$) and rotational ($\bm\Omega$) parts obtained from the
velocity gradient projected into that frame; the $2\times2$ closed form is in
\texttt{logconf\_local\_rhs}. The local update has been verified to machine
precision against the analytic step-strain relaxation, the steady-shear viscometric
functions (including the first normal-stress difference $N_1=2G\tau^2\dot\gamma^2$),
and the $\tau\!\to\!\infty$ elastic limit (Section~\ref{sec:vv}).

\section{Spatial discretization: the MAC staggered grid}
On an $N_x\times N_y$ grid with spacing $(\Delta x,\Delta y)$:
\begin{center}
\begin{tabular}{lll}
\toprule
field & location & array shape \\
\midrule
pressure $p$ & cell centres & $(N_y,N_x)$ \\
$u$ (x-velocity) & $x$-faces & $(N_y,N_x{+}1)$ \\
$v$ (y-velocity) & $y$-faces & $(N_y{+}1,N_x)$ \\
$\xiv,\phi,\bm\psi$ & cell centres & $(N_y,N_x)$ \\
\bottomrule
\end{tabular}
\end{center}
The discrete divergence (\texttt{divergence}) maps face velocities to cell centres
and the pressure gradient (\texttt{gradient\_p\_u/v}) maps centres back to faces;
they are exact transposes, $\mathrm{div}=-\mathrm{grad}^{\!\top}$. Consequently the
Laplacian $\mathrm{div}\,\mathrm{grad}$ used in the pressure solve is consistent with
the operators used in the correction, which is what makes the projection exact.

\section{Pressure projection and Poisson solvers}
A fractional-step (Chorin) projection is used. Given a provisional velocity
$\uv^{*}$ the pressure (here $\phi=p\,\Delta t/\rho$) solves
\begin{equation}
\nabla^2\phi = \frac{\rho}{\Delta t}\,\grad\!\cdot\!\uv^{*},
\qquad
\uv^{n+1} = \uv^{*}-\frac{\Delta t}{\rho}\grad\phi .
\end{equation}
Two direct $O(N\log N)$ solvers are provided:
\begin{itemize}[nosep]
  \item \textbf{Neumann walls} (\texttt{poisson\_eigs\_neumann},
  \texttt{solve\_poisson\_neumann}): a discrete cosine transform (DCT-II); the
  homogeneous Neumann pressure condition means the projection leaves the wall-normal
  velocity unchanged, which is exploited to impose prescribed inflow/outflow
  boundary velocities (e.g.\ uniform extension) as long as the net wall flux is zero.
  \item \textbf{Periodic} (\texttt{poisson\_eigs\_periodic},
  \texttt{solve\_poisson\_periodic}): an FFT.
\end{itemize}
Because the operators are consistent, the resulting velocity is divergence-free to
$\sim\!10^{-14}$ regardless of the flow (verified in the benchmarks). The collocated
solver additionally offers a PyAMG multigrid solve for variable-coefficient
(variable-density) projections.

\section{Time integration}
The default scheme is an explicit predictor followed by the projection:
\begin{enumerate}[nosep]
  \item Advect the reference map (and $\bm\psi$, $\bten_e$) with the cell-centred
  velocity (Section~\ref{sec:advect}), then mask to the solid and extrapolate.
  \item Compute the blended stress and assemble the face force $\bm f$.
  \item Predictor (\texttt{momentum\_predictor}, central advection $+$ diffusion,
  forward Euler) producing $\uv^{*}$, with optional face body forces $f_u,f_v$ added
  as $f/\rho$.
  \item Project to obtain the divergence-free $\uv^{n+1}$ and $p$.
\end{enumerate}
Three momentum predictors are provided: lid-driven (top wall moves at $U_{\text{lid}}$,
others no-slip), free-slip box, and periodic. The collocated solver additionally
offers an RK4 momentum step (\texttt{momentum\_step\_rk4}).
The explicit time step obeys the advective, viscous, and elastic-wave CFL limits,
\begin{equation}
\Delta t \le \min\!\bigl(\underbrace{C\,\Delta x/U}_{\text{advection}},\;
\underbrace{C_\nu\,\Delta x^2/\nu}_{\text{viscous}},\;
\underbrace{C\,\Delta x/c_s}_{\text{elastic wave}}\bigr),
\qquad c_s=\sqrt{\mu_s/\rho},
\label{eq:cfl}
\end{equation}
and, for viscoelasticity, additionally $\Delta t\lesssim\tau$ from the stiff
relaxation. The last three are lifted by the implicit-explicit (IMEX) stack below.

\subsection{Implicit--explicit (IMEX) strategy stack}
\label{sec:imex}
The stiff terms in \eqref{eq:cfl} are treated implicitly one at a time, forming a
\emph{stack} of strategies (each contains the previous one). Advection stays explicit
throughout (the semi-Lagrangian backtrace of Section~\ref{sec:advect} is already
unconditionally stable).

\paragraph{Implicit viscosity (removes $C_\nu\Delta x^2/\nu$).}
The viscous term is advanced by backward Euler, giving a Helmholtz solve for the
predictor,
\begin{equation}
\bigl(I-\Delta t\,\nu\nabla^2\bigr)\,\uv^{*}
   = \uv^{n} - \Delta t\,(\uv\!\cdot\!\nabla)\uv^{n} + \Delta t\,\bm f/\rho .
\label{eq:helmholtz}
\end{equation}
On periodic (and free-slip / Neumann) directions the Helmholtz operator is
diagonalized by the \emph{same} FFT/DCT used for the pressure Poisson solve, so
\eqref{eq:helmholtz} costs one extra transform per component
(\texttt{solve\_helmholtz\_periodic}, \texttt{lap\_eigs\_periodic}). For the
Dirichlet (no-slip $+$ moving lid) cavity the operator is not diagonalized by the DCT,
so it is solved matrix-free by conjugate gradients reusing the exact ghost-cell viscous
stencil (\texttt{momentum\_predictor\_lid\_imex}, \texttt{\_cg\_helmholtz}); the moving
lid enters as a right-hand-side inhomogeneity, and the operator generalizes to the
variable one-fluid viscosity.

\paragraph{Exact/exponential relaxation (removes $\Delta t\lesssim\tau$).}
The log-conformation update is split (Strang) into
$\text{relax}(\Delta t/2)\!\to\!\text{stretch}(\Delta t)\!\to\!\text{relax}(\Delta t/2)$,
with the relaxation half-steps integrated \emph{analytically},
\begin{equation}
\bten_e \leftarrow I + (\bten_e-I)\,e^{-\Delta t/2\tau},
\label{eq:relaxexact}
\end{equation}
the closed-form solution of $\dot{\bten}_e=-\tfrac1\tau(\bten_e-I)$. This is
$A$-stable for any $\tau$ and SPD-preserving (a convex blend of two SPD tensors), and
for $\tau\to\infty$ reduces exactly to the explicit stretch step
(\texttt{logconf\_local\_step\_strang}, \texttt{relax\_exact}).

\paragraph{Linearly-implicit elastic stabilizer (removes $C\,\Delta x/c_s$).}
The nonlinear elastic force $\nabla\!\cdot\!\bm\sigma_{\!e}$ is kept explicit
(physics), but an $O(\Delta t^2)$ wave operator $\Delta t^2 c_s^2\,\chi\,\nabla^2$
($\chi$ the solid indicator) is added \emph{implicitly} to \eqref{eq:helmholtz}. A
1-D von-Neumann analysis of the linear elastic wave
$\uv_t=c_s^2 d_{xx},\ d_t=\uv$ gives the amplification
$|\lambda|=1/\sqrt{1+\Delta t^2 c_s^2 k^2}\le1$, i.e.\ unconditional stability with
mild damping --- \emph{provided the elastic force is placed inside the implicit
right-hand side} (adding it after the solve is unstable at large $\Delta t$). A
frozen-coefficient split keeps a single transform solve
(\texttt{momentum\_predictor\_periodic\_imex\_elastic}). This is stage~1: the
isotropic-Laplacian stabilizer adds numerical damping at large $\Delta t$; the
fully-consistent Newton/JFNK tangent that removes it is future work
(Section~\ref{sec:limits}).

\paragraph{Unified selector.} \texttt{resolve\_integrator} maps a single
\texttt{integrator} keyword to the internal flags, so every MAC driver exposes one
knob,
\[
\texttt{explicit}\ \subset\ \texttt{imex}\ \subset\ \texttt{imex-elastic},
\]
with the legacy booleans (\texttt{implicit\_visc}, \texttt{imex},
\texttt{elastic\_imex}) retained as aliases. Each level makes one more stiff term of
\eqref{eq:cfl} implicit.

\begin{center}\small
\begin{tabular}{@{}llll@{}}
\hline
\texttt{integrator} & implicit terms & lifts & drivers\\
\hline
\texttt{explicit}     & --- & (baseline) & all\\
\texttt{imex}         & viscosity $+$ relaxation & $\Delta x^2/\nu$, $\tau$ & all\\
\texttt{imex-elastic} & $+$ elastic wave & $+\,\Delta x/c_s$ & viscoelastic\\
\hline
\end{tabular}
\end{center}

\subsection{Reference-map advection schemes}
\label{sec:advect}
The dispatcher \texttt{advect\_reference\_map(\dots, scheme)} selects among
(\texttt{semilagrangian} bilinear RK4 --- default and most robust;
\texttt{semilagrangian\_cubic} monotone Catmull--Rom;
\texttt{central2} $+$ SSP-RK3;
\texttt{weno5} $+$ SSP-RK3;
\texttt{conservative} Jain Eq.~26 with the divergence-free face velocity). The choice
trades robustness against formal order; see Section~\ref{sec:conv}.

\section{Interface, extrapolation and reinitialization}
\label{sec:interface}
\paragraph{Smoothed indicator.} The blend \eqref{eq:blend} uses a smoothed Heaviside
$H=\tfrac12\bigl(1+\tanh(\phi/w_t)\bigr)$ (\texttt{smoothed\_heaviside}) over a
transition width $w_t\sim 2\Delta x$, giving a \emph{diffuse} interface of a few
cells.
\paragraph{Extrapolation.} After advection the reference map is defined only inside
the solid; \texttt{extrapolate\_reference\_map} fills a narrow band outside the zero
contour (a few layers) so that gradients and the rebuilt level set are well defined.
\paragraph{Level-set rebuild.} The interface is reconstructed each step from the
reference map, $\phi=\phi_0(\xiv)$ (\texttt{rebuild\_phi\_from\_reference\_map}), or
by bilinear interpolation of a stored grid level set.
\paragraph{Reinitialization.} Two notions are distinguished:
\begin{itemize}[nosep]
  \item \emph{Level-set} reinitialization to a signed distance (FMM via
  \texttt{reinitialize\_phi\_fmm}, requires \texttt{scikit-fmm}; or the PDE scheme).
  Per-step FMM is not volume-preserving; it is applied infrequently.
  \item \emph{Reference-map} reinitialization (Rycroft F-storage): when the
  per-epoch distortion grows large, reset $\xiv=\xv$, fold the current deformation
  into a stored gradient $\Ften_{\text{stored}}\!\leftarrow\!\Ften\Ften_{\text{stored}}$,
  and adopt the current shape as the new reference. Stress then uses
  $\Ften_{\text{tot}}=\Ften\Ften_{\text{stored}}$. This lets a solid under sustained
  \emph{smooth} shear run indefinitely. It does \emph{not} cure folds caused by
  rigid rotation winding (use the viscoelastic decoupling, Section~\ref{sec:ve}) or
  by under-resolved contact crushing (use resolution / lubrication).
\end{itemize}

\section{Surface tension}
Surface tension is added as a balanced-force continuum surface force (CSF) evaluated
\emph{on the faces}:
\begin{equation}
\bm f_{\text{st}} = -\gamma\,\kappa\,\grad H ,
\end{equation}
with curvature $\kappa=-\grad\!\cdot\!(\grad\phi/|\grad\phi|)$
(\texttt{compute\_curvature}) and the \emph{same compact face gradient} used for the
pressure (\texttt{interfacial\_force\_faces}). Using the identical operator for the
force and the pressure gradient makes the discrete balance exact for a static drop,
which suppresses parasitic currents (measured $\sim\!5.5\times$ smaller than the
collocated discretization). The static Laplace jump is reproduced to $\sim\!0.4\%$.

\section{Contact between solids}
A repulsive \emph{body} force between near-contacting solids is curl-free and is
therefore removed by the exact projection (it has no effect). Contact is instead
added as a trace-free \emph{stress} (Rycroft \emph{et al.}, 2018,
Eqs.~4.10--4.12), implemented in \texttt{contact\_stress}:
\begin{equation}
\bm\tau_{\text{col}} = -\eta\,\min\!\bigl(f(\phi_i),f(\phi_j)\bigr)\,(G_i{+}G_j)
\Bigl(\bm n\!\otimes\!\bm n-\tfrac12\bm I\Bigr),
\end{equation}
active where the two diffuse bands overlap (width $\varepsilon\sim 3\Delta x$), with
$\bm n$ from the level-set gradient. Its divergence is a momentum-conserving force
that survives the projection; the stiffness $\eta$ controls the rebound.
Solid--wall contact uses the same routine against a wall level set
$\phi_{\text{wall}}=\min(x,1{-}x,y,1{-}y)$.
\paragraph{Limitation \& outlook.} The penalty contact prevents overlap but is not a
quantitative model of the thin lubricating film; near-contact at gaps below the grid
scale is under-resolved (it folds unless refined). A \emph{lubrication} reformulation
(adapting Fai \& Rycroft, 2018, to the diffuse-band RMT, using the signed-distance
fields as a height function) is planned to make the contact force physical and
resolution-robust.

\section{Verification and validation}
\label{sec:vv}
\subsection{Operators and projection}
The discrete divergence/gradient transpose property and the projection are checked
in \texttt{tests/test\_mac.py}; the divergence of the projected field is
$\sim\!10^{-14}$.

\subsection{Spatial convergence}
\label{sec:conv}
\begin{center}
\begin{tabular}{lll}
\toprule
case & quantity & observed order \\
\midrule
Taylor--Green (pure fluid) & $L_2$ velocity & $2.01\!\to\!2.06\!\to\!2.25$ \\
Disc-in-Taylor--Green (FSI), bilinear SL & velocity & $1.31$ \\
\quad cubic SL & velocity & $1.72$ \\
\quad central2 & velocity & $1.74$ \\
\quad WENO5 & velocity & $2.10$ \\
\bottomrule
\end{tabular}
\end{center}
The pure-fluid solver is second order. The FSI velocity order is below~2 and is
\emph{capped by the diffuse interface} (the smeared stress jump), not by the
advection scheme: improving the scheme raises the order only up to the
diffuse-interface ceiling (issue tracked as ``2nd-order fields require a sharper
interface''). The divergence stays at machine zero across all schemes and
resolutions.

\subsection{Benchmarks against the literature}
\begin{itemize}[nosep]
  \item \textbf{Lid-driven cavity} vs.\ Ghia \emph{et al.} (1982): RMS error in
  $u(y)$ at $x{=}0.5$ of $1.7\times10^{-3}$ (Re~100) and $2.8\times10^{-2}$
  (Re~1000), $N{=}129$.
  \item \textbf{Soft disc in lid-driven cavity} vs.\ Sugiyama \emph{et al.} (2011)
  / Jain \emph{et al.} (2019): centroid trajectory matches; the deformed interface
  is grid-converged ($N{=}64\approx N{=}128$).
  \item \textbf{Surface tension}: static Laplace jump to $0.4\%$; a square relaxes
  to a circle (circularity $0.999$, area conserved to $\sim2\%$ with infrequent
  reinitialization).
  \item \textbf{Contact}: two-disc collision rebound (monotone in $\eta$);
  multi-body and multi-shape stirring.
\end{itemize}

\subsection{Viscoelasticity}
The constitutive update (both the plain $\bten_e$ and the log-conformation form) is
verified to machine precision against analytic responses: step-strain relaxation
$\sigma_{xy}=G\gamma e^{-t/\tau}$ ($\sim\!10^{-10}$), steady shear
$\sigma_{xy}=G\tau\dot\gamma$ and $N_1=2G\tau^2\dot\gamma^2$ ($\sim\!10^{-12}$), and
the elastic limit $\bten_e=\Ften\Ften^{\!\top}$ ($\sim\!10^{-13}$); see
\texttt{tests/test\_viscoelastic.py}. Under imposed uniform planar extension the
solver reproduces the analytic plateau $b_e^{xx}\!\to\!1/(1-2\,\mathrm{Wi})$ to four
digits ($\tau{=}0.25\!\to\!1.3333$, $\tau{=}0.5\!\to\!1.997$ vs.\ $2.0$) and is
resolution-independent ($N{=}64/96/128$ overlap), while the elastic case diverges as
$e^{2\varepsilon t}$ (Fig.~\texttt{viscoelastic\_uniform\_money.png}).

\subsection{IMEX time integration}
The strategy stack of Section~\ref{sec:imex} is verified to reproduce the explicit
results before claiming any lifted step. \emph{Implicit viscosity:} on Taylor--Green
the scheme is stable and accurate to $\Delta t=50\times$ the explicit viscous limit
(where explicit diverges) and keeps second-order spatial accuracy. \emph{Exact
relaxation:} step-strain is reproduced to machine precision at \emph{any} $\Delta t$,
whereas the explicit step errs by $0.29$ at $\Delta t/\tau=2$. \emph{Combined IMEX}
(four-roll blob) reproduces the explicit $b_e^{xx}$ to $0.55\%$ at a matched step and
completes within ${\sim}2\%$ at $2\!-\!4\times$ the explicit viscous CFL where explicit
diverges. \emph{Wall solver:} the lid cavity reproduces the Ghia steady state and the
soft disc reproduces the (Sugiyama-validated) centroid trajectory to $2\times10^{-3}$.
\emph{Elastic stabilizer (stage~1):} consistent ($0.44\%$ vs.\ explicit at $0.1\times$
the elastic CFL) and stable to $8\times$ the elastic CFL where the explicit elastic
force diverges; see \texttt{tests/test\_\{mac,ve\_imex\_fsi,imex\_wall,elastic\_imex\}.py}.

\section{Limitations}
\label{sec:limits}
\begin{itemize}[nosep]
  \item Two-dimensional; single density by default (variable density only in the
  collocated AMG path).
  \item Diffuse-interface FSI fields converge below second order
  (Section~\ref{sec:conv}); a sharp-interface treatment is future work.
  \item Penalty contact is non-quantitative near contact (lubrication planned).
  \item The IMEX stack (Section~\ref{sec:imex}) lifts the viscous, relaxation and
  (stage~1) elastic-wave limits; the elastic stabilizer is only linearly-implicit and
  adds numerical damping at large $\Delta t$. The fully-consistent Newton/JFNK elastic
  kernel (removing the damping, giving accurate large-$\Delta t$ steps) is future work.
  \item Long, high-resolution runs are memory bound; frame data should be streamed
  to disk (the showcase driver does this).
\end{itemize}

\section{Module map}
\begin{center}
\begin{tabular}{ll}
\toprule
file & contents \\
\midrule
\texttt{pyRMT/mac.py} & staggered grid, exact projection, momentum predictors, \\
 & periodic/Neumann Poisson, IMEX Helmholtz (FFT \& CG), \\
 & \texttt{resolve\_integrator} selector, interfacial \& contact stress \\
\texttt{pyRMT/functions.py} & collocated solver, advection schemes, neo-Hookean stress, \\
 & extrapolation, reinitialization, AMG projection \\
\texttt{pyRMT/viscoelastic.py} & UCM $\bten_e$, log-conformation, $\exp/\log$ of $2{\times}2$ SPD \\
\texttt{pyRMT/interpolators.py} & bilinear / bicubic interpolation (Numba) \\
\texttt{pyRMT/utils.py} & finite-difference gradient/Laplacian helpers \\
\texttt{pyRMT/output.py} & kinetic / strain energy, dissipation, HDF5 output \\
\bottomrule
\end{tabular}
\end{center}

\section*{References}
\begin{enumerate}[nosep]
  \item K.~Kamrin, C.~H.~Rycroft, J.-C.~Nave, \emph{Reference map technique for
  finite-strain elasticity and fluid--solid interaction}, JMPS (2012).
  \item C.~H.~Rycroft \emph{et al.}, \emph{Reference map technique for
  incompressible fluid--structure interaction}, JFM (2020).
  \item S.~S.~Jain, K.~Kamrin, A.~Mani, \emph{A conservative and non-dissipative
  Eulerian formulation for the simulation of soft solids in fluids}, JCP (2019).
  \item R.~Fattal, R.~Kupferman, \emph{Constitutive laws for the matrix-logarithm of
  the conformation tensor}, JNNFM (2004).
  \item T.~G.~Fai, C.~H.~Rycroft, \emph{Lubricated immersed boundary method in two
  dimensions}, JCP (2018).
  \item U.~Ghia, K.~N.~Ghia, C.~T.~Shin, \emph{High-Re solutions \dots}, JCP (1982).
\end{enumerate}

\end{document}
